\documentclass[11pt]{article}
\usepackage[a4paper,margin=1in]{geometry}
\usepackage{fontspec}
\usepackage[boldfont=false]{xeCJK}
\setCJKmainfont{FandolSong-Regular.otf}
\usepackage{amsmath}
\usepackage{amssymb}
\usepackage{array}
\usepackage{booktabs}
\usepackage{graphicx}
\usepackage{adjustbox}
\usepackage{enumitem}
\setlist[itemize]{topsep=2pt,partopsep=0pt,parsep=0pt,itemsep=2pt,leftmargin=1.4em}
\setlist[enumerate]{topsep=2pt,partopsep=0pt,parsep=0pt,itemsep=2pt,leftmargin=1.6em}
\usepackage{parskip}
\usepackage{fvextra}
\fvset{breaklines=true,breakanywhere=true,fontsize=\small}
\setlength{\parindent}{0pt}

\begin{document}

\begin{center}
  {\LARGE\bfseries Transport in animals}\\[0.35em]
  {\large Igcse Biology}
\end{center}
\vspace{0.6em}

\section*{The circulatory system}

The \textbf{circulatory system} 循环系统 carries blood around the body. It has three parts: a \textbf{pump} 泵 (the \textbf{heart} 心脏), a network of \textbf{blood vessels} 血管, and \textbf{valves} 瓣膜 that keep the blood flowing one way only.

\subsection*{Single and double circulation (Supplement)}

\begin{itemize}
\item A fish has a \textbf{single circulation} 单循环: the blood passes through the heart \textbf{once} on each trip around the body.

\item A \textbf{mammal} 哺乳动物 has a \textbf{double circulation} 双循环: the blood passes through the heart \textbf{twice} on each full trip — once on the way to the lungs, and once on the way to the rest of the body.

\item The advantage of a double circulation: the blood can be pumped again at \textbf{high pressure} 压力 before going to the body, so it travels faster and delivers \textbf{oxygen} 氧气 quickly.

\end{itemize}

\section*{The heart}

The heart is made of \textbf{muscle} 肌肉. A wall called the \textbf{septum} 隔膜 divides it into a left side and a right side. On each side there is an upper chamber, the \textbf{atrium} 心房, and a lower chamber, the \textbf{ventricle} 心室. Valves between the chambers stop the blood flowing backwards. The heart muscle is fed with blood by the \textbf{coronary arteries} 冠状动脉.

\textbf{(Supplement)} The wall of the left ventricle is thicker than the right, because it must pump blood all the way around the body; the right side only pumps blood to the nearby lungs. The atria have thin walls, as they only push blood into the ventricles just below them.

The septum keeps \textbf{oxygenated blood} 含氧血 on the left side apart from \textbf{deoxygenated blood} 缺氧血 on the right side, so the two never mix.

\textbf{(Supplement)} The valves between each atrium and ventricle are the \textbf{atrioventricular valves} 房室瓣. The valves at the exit of each ventricle are the \textbf{semilunar valves} 半月瓣.

\subsection*{How the heart beats (Supplement)}

\begin{enumerate}
\item The atria \textbf{contract} 收缩 and push blood down into the ventricles.

\item The ventricles contract and force blood out into the arteries; the valves snap shut so blood cannot flow back.

\item The muscle relaxes and the heart fills with blood again.

\end{enumerate}

Blood is pumped away from the heart in arteries and returns to the heart in veins.

\subsection*{Heart rate and exercise}

You can check the heart's activity with an \textbf{ECG}, by feeling the \textbf{pulse} 脉搏 in an artery, or by listening to the valves closing. During exercise the \textbf{heart rate} 心率 goes up, so blood reaches the muscles faster, bringing the extra oxygen and glucose they need.

\subsection*{Coronary heart disease}

\textbf{Coronary heart disease} 冠心病 happens when the coronary arteries become narrow or blocked, so the heart muscle cannot get enough oxygen. \textbf{Risk factors} 风险因素 include a diet high in fat, lack of exercise, stress, smoking, family history (genetics), older age, and being male. A healthy diet and regular exercise lower the risk.

\section*{Blood vessels}

There are three kinds of blood vessel. The space inside a vessel is its \textbf{lumen} 管腔.

\begin{center}
\adjustbox{max width=\linewidth}{%
\begin{tabular}{lllll}
\toprule
\textbf{Vessel} & \textbf{Wall} & \textbf{Lumen} & \textbf{Valves?} & \textbf{Job} \\
\midrule
\textbf{arteries} 动脉 & thick, muscular, elastic & narrow & no & carry blood \textbf{away} from the heart, at high pressure \\
\textbf{veins} 静脉 & thin & wide & yes & carry blood \textbf{back} to the heart, at low pressure \\
\textbf{capillaries} 毛细血管 & one cell thick & very narrow & no & let \textbf{oxygen}, glucose and wastes pass between blood and cells \\
\bottomrule
\end{tabular}}
\end{center}

\textbf{(Supplement)} Arteries have thick elastic walls to cope with the high pressure from the heart. Veins have valves because their blood is at low pressure and could otherwise flow backwards. Capillaries are very thin (one cell thick) and very narrow, giving a large \textbf{surface area} 表面积 and a short distance for fast exchange.

\subsection*{Main blood vessels}

\begin{center}
\adjustbox{max width=\linewidth}{%
\begin{tabular}{lll}
\toprule
\textbf{Connects} & \textbf{Artery (carries blood out)} & \textbf{Vein (carries blood in)} \\
\midrule
heart ↔ body & \textbf{aorta} 主动脉 & \textbf{vena cava} 腔静脉 \\
heart ↔ lungs & \textbf{pulmonary artery} 肺动脉 & \textbf{pulmonary vein} 肺静脉 \\
↔ kidney & \textbf{renal artery} 肾动脉 & \textbf{renal vein} 肾静脉 \\
\bottomrule
\end{tabular}}
\end{center}

\textbf{(Supplement)} The liver is served by the \textbf{hepatic artery} 肝动脉, the \textbf{hepatic vein} 肝静脉, and the \textbf{hepatic portal vein} 肝门静脉, which brings blood from the gut to the liver.

\section*{Blood}

Blood has four parts:

\begin{center}
\adjustbox{max width=\linewidth}{%
\begin{tabular}{ll}
\toprule
\textbf{Part} & \textbf{Function} \\
\midrule
\textbf{red blood cells} 红细胞 & carry oxygen, using the red \textbf{pigment} 色素 \textbf{haemoglobin} 血红蛋白 \\
\textbf{white blood cells} 白细胞 & defend the body against disease \\
\textbf{platelets} 血小板 & help the blood to \textbf{clot} 凝血 \\
\textbf{plasma} 血浆 & a liquid that carries blood cells, \textbf{ions} 离子, \textbf{nutrients} 营养物质, \textbf{urea} 尿素, \textbf{hormones} 激素 and \textbf{carbon dioxide} 二氧化碳 \\
\bottomrule
\end{tabular}}
\end{center}

\subsection*{Defending the body, and clotting}

White blood cells fight \textbf{pathogens} 病原体 (the microbes that cause disease) in two ways:

\begin{itemize}
\item \textbf{phagocytes} 吞噬细胞 carry out \textbf{phagocytosis} 吞噬作用 — they surround and digest the pathogens.

\item \textbf{lymphocytes} 淋巴细胞 make \textbf{antibodies} 抗体, which stick to the pathogens and destroy them.

\end{itemize}

When you cut yourself, the blood \textbf{clots} to seal the wound. This stops blood loss and stops pathogens getting in. \textbf{(Supplement)} During clotting, a soluble protein called \textbf{fibrinogen} 纤维蛋白原 is changed into threads of \textbf{fibrin} 纤维蛋白. The threads form a net that traps blood cells and makes a solid clot.

\section*{Exam tips}

\begin{itemize}
\item Circulatory system = pump (heart) + vessels + valves for one-way flow.

\item Double circulation (mammals): blood passes through the heart twice. The left ventricle wall is thickest because it pumps to the whole body.

\item Arteries carry blood \textbf{away} (thick, high pressure); veins carry it \textbf{back} (valves, low pressure); capillaries are for exchange (one cell thick).

\item Blood = red cells (oxygen, haemoglobin) + white cells (defence) + platelets (clotting) + plasma (transport).

\item Phagocytes engulf pathogens; lymphocytes make antibodies.

\end{itemize}


\end{document}
